\documentclass[10pt,a4paper]{amsart}
\usepackage[margin=19mm]{geometry}
\usepackage{amsmath,amssymb,mathtools}
\usepackage[scr=rsfs,cal=euler]{mathalfa}
\usepackage{xfrac}
\usepackage[unicode,hidelinks,bookmarks=false]{hyperref}
\newcommand{\ie}{i.e.\ }
\newcommand{\cf}{cf.\ }
\newcommand{\resp}{respectively\ }
\newcommand{\Subsection}{\subsection}
\providecommand{\nobreakdash}{\nobreak\hbox{-}}
\setlength{\emergencystretch}{2em}
\multlinegap0pt
\usepackage{mdframed}
\usepackage{mdframed}
\usepackage{mdframed}
\usepackage{mdframed}
\usepackage{amssymb}
\usepackage{mathtools}
\usepackage{xcolor}
\usepackage{mdframed}
\usepackage[shortlabels]{enumitem}
\newtheorem{proposition}{Proposition}
\title{$sp$-Homogeneous Linear Orderings}
\author{Wesley Calvert, Douglas Cenzer, David Gonzalez, Valentina Harizanov, Keng Meng Ng}
\date{Source: arXiv:2509.25005v2 (CC BY 4.0)}
\begin{document}
\maketitle

\noindent\emph{Source and licence.} $sp$-Homogeneous Linear Orderings. Version arXiv:2509.25005v2 (CC BY 4.0), distributed by arXiv under the Creative Commons Attribution 4.0 International licence, http://creativecommons.org/licenses/by/4.0/, as recorded in the arXiv licence metadata for this exact version and shown on the abstract page https://arxiv.org/abs/2509.25005v2. Full licence notice: \texttt{licenses/arxiv-2509.25005-CC-BY-4.0.txt}.\par\smallskip

\noindent\textit{Editorial scope.} Counted unit: the article's proposition Le1 (source0.tex lines 650--658). The article's complete original proof of it is delivered with it (source0.tex lines 660--728, one \texttt{proof} environment of 69 lines). No mathematics is written, repaired or re-derived: the statement and the proof are the article's own lines, in the article's own notation, and no retained line is substituted or re-flowed.  Retained uncounted as the source-local context the proof needs: lines 556--558.  Nothing was omitted from any retained span, so the package carries no omission record.  No retained line contains a citation, so the wrapper carries no bibliography and no bibliography entry.  No retained line contains a figure environment, an image-inclusion command or drawing code, so no image is delivered and the package declares no asset dependency.  Editorial formatting only, in addition: an AMS \texttt{amsart} preamble that restates the article's own declarations and macros used by the retained lines, namely the environments \texttt{proposition} on separate counters, together with the standard packages those lines need.  The retained environments print in this document as Proposition 1 in order of appearance, whereas the article prints the same items with the numbering of its own full document.  The complete native source of the article as distributed in the arXiv e-print for this exact version is bundled unchanged as \texttt{sources/185989/source0.tex}.
\begin{proposition} \label{shnotisom} 
There are two computable copies of $Sh(\omega + 1)$ which are not $\Delta^0_3$ isomorphic.
\end{proposition}

\begin{proposition} \label{Le1} Let $n>1$ be finite.
\begin{enumerate}
\item $LI$ is a $\Pi^0_1$ set.
\item $\{e: L_e \cong \eta\}$ is $\Pi^0_2$ complete.
\item  $\{e: L_e \cong n \cdot \eta\}$ is $\Pi^0_3$ complete.
\item $\{e \in LI: L_e \cong \omega \cdot \eta \}$,  $\{e \in LI: L_e \cong \omega^* \cdot \eta \}$, and $\{e \in LI: L_e \cong \zeta \cdot \eta\}$ are all $\Pi^0_4$ complete.
\item $\{e \in LI: L_e \cong \omega\}$,  $\{e \in LI: L_e \cong \omega^*\}$, and $\{e \in LI: L_e \cong \zeta\}$ are all $\Pi^0_3$ complete.
\end{enumerate}
\end{proposition}

\begin{proof} The upper bounds on the complexity are easy to see. 
We will just show that  $\{e: L_e \cong n \cdot \eta\}$ is $\Pi^0_3$.
That is, for $e \in LI$, $L_e \cong n \cdot \eta$ if and only if every block size is $n$, and 
\begin{itemize}
\item for all $x$, $|[x]| = n$;
\item for any $x$, there exist $y_1 < y_2 < y_3 < x$;
\item for any $x$, there exist $y_1 > y_2 > y_3 > x$;
\item for any $x_0 < x_1$, there exist $y_1,y_2,y_3$ such that 
$x < y_1 < y_2 < y_3 < x_1$. 
\end{itemize}

\medskip

For the completeness:

(2) We obtain a reduction $f$ from the $\Pi^0_2$ complete set $Inf = \{e: W_e\ \text{is infinite}\}$ so that $L_{f(e)} \cong \eta$ if $e \in Inf$,and $L_{f(e)} \cong \zeta$ if $e \notin Inf$. Initially, $L = L_{f(e)} = \{0\}$. After stage $s$, we will have a finite ordering $L^s = a_0^s < a_1^s < \dotsb < a_{n_s -1}$ of size $n_s$ with domain $n_s$.  At stage $s+1$, there are two cases.

\begin{itemize}
    \item If no new element is enumerated into $W_e$, then we add $n_s$ on the left of $L^s$,  and we add $n_s+1$ on the right. Thus $n_{s+1} = n_s+2$.
    \item If a new element is enumerated into $W_e$, then we add $n_s$ on the left of $L^s$,
    we add $n_s+i+1$ between $a_i^s$ and $a_{i+1}^s$, and we add $2n_s$ on the right. Thus $n_{s+1} = 2n_s+1$. 
\end{itemize}

We have that $L = L_{f(e)}$ is the union of the orders $L^s$. If $W_e$ is finite, then clearly $L_{f(e)} \cong \zeta$. If $W_e$ is infinite, then the construction ensures that $L$ is dense and without endpoints, so that $L_{f(e)} \cong \eta$. 


(3) We use the $\Pi^0_2$ completeness of $Inf$. Given an arbitrary $\Pi^0_3$ predicate $Q(e) \equiv (\forall m) P(e,m)$, where $P$ is $\Sigma^0_2$, there is a c.e.\ relation $R$ such that $P(e,m)$ if and only if $\{x: R(e,m,x)\}$ is finite. Here we may assume that if $\neg Q(e)$, then $\neg P(e,m)$ for co-finitely many $m$. 

It suffices to define a computable function $\phi$ so that, for all $e$, $Q(e)$ if and only if $L_{\phi(e)} \cong n \cdot \eta$.

Then $L_{\phi(e)}$ is defined as the sum $L_0 + n + L_1 + n +  \dotsb$, where for each $m$, $L_m \cong \eta$ if $P(e,m)$ and $L_m \cong n\cdot \eta$ if $\neg P(e,m)$. Thus, if $Q(e)$, then each $L_m \cong n \cdot \eta$, so that $L_{\phi(e)} \cong n \cdot \eta$, and otherwise at least one of the $L_m \cong \eta$, so that $L_{\phi(e)}$ is not isomorphic to $n \cdot \eta$. 

For the construction of the component $L_m$, we are building a copy of $n \cdot \eta$, but every time a new element $x$ satisfies $R(e,m,x)$, we break up our size $n$ blocks $y_1 < y_2 < \dots < y_n$ into $n$ blocks of size $n$ by starting to build a copy of $n \cdot \eta$ between each successor pair and adding new elements to expand each $y_i$ to a block of size $n$. Thus, if $P(e,m)$, then this expansion only happens finitely often, so that $L_m \cong n \cdot \eta$. 
If $\neg P(e,m)$, then the expansion happens infinitely often, so that no elements of $L_m$ will have (permanent) successors, and thus $L_m \cong \eta$, so that $L \cong n \cdot \eta + n + n \cdot \eta + \dotsb + n + n \cdot \eta + n + \eta + n + \eta + \cdots\not\cong n\cdot\eta$. 


\medskip

(4) Let $P(x) \equiv (\forall a)R(a,x)$, where $R$ is $\Sigma^0_3$. We use the technique from the proof of Proposition \ref{shnotisom} to build $L_{f(x)}$ as a dense family of blocks. Each block will have the form $M_0 + M_1 + \cdots$, where $M_a$ is used to check whether $R(a,x)$. The component $M_a$ will have finite size if $R(a,x)$, and it will have type $\omega$ if $\neg R(a,x)$. As in the proof of Proposition \ref{shnotisom}, we use the $\Sigma^0_3$ completeness of $Cof$ to define $M_a$ so that:

\begin{itemize}
\item 
 When $R(a,x)$, $M_a$ will have finite size $n_e$, where $n_e$ is the least $n$ such that all elements $\geq n$ belong to the appropriate c.e.\ set, and 

\item When $\neg R(a,x)$, $M_a$ will have type $\omega$. 
\end{itemize}

The discarded elements outlined in the proof of Proposition \ref{shnotisom} are placed individually in new blocks $M_i$ that have not yet been started. At each stage of the construction, in addition to updating the $M_a$, we start to build new blocks between any two blocks and add new blocks at each end, so that we will get a dense family of blocks. 

If $P(x)$, then each $M_a$ will be finite, so that $M_0 + M_1 + \cdots$ will have type $\omega$ and therefore $L_{f(x)} \cong \omega \cdot \eta$. 

If $\neg P(x)$, then eventually each $M_a$ will have type $\omega$,  so that $M_0 + M_1 + \cdots$ will have type $\omega \cdot \omega$ and therefore $L_{f(x)} \cong \omega^2 \cdot \eta$. 

The proof is very similar when considering $\omega^*\cdot\eta$ or $\zeta\cdot\eta$.

\medskip

(5) We again use the $\Pi^0_2$ completeness of $Inf$. Given an arbitrary $\Pi^0_3$ predicate $Q(e) \equiv (\forall m) P(e,m)$, where $P$ is $\Sigma^0_2$, there is a c.e.\ relation $R$ such that $P(e,m)$ if and only if $\{x: R(e,m,x)\}$ is finite. Here we may assume that if $\neg Q(e)$, then $\neg P(e,m)$ for co-finitely many $m$. Also, we may assume that at any stage $s+1$, there is exactly one pair $(m,x)$ such that $e,m,x$ comes into $R$ at stage $s+1$. 

It suffices to define a computable function $\phi$ so that, for all $e$, $Q(e)$ if and only if $L_{\phi(e)} \cong \omega$. 

The ordering $L_{\phi(e)}$ is defined as the sum $L_0 + 3 + L_1 + 3 +  \dotsb$, where for each $m$, $L_m \cong \omega$, if $\neg P(e,m)$, and $L_m$ is finite and has size $k_m$ if $\{x: P(e,m,x)\}$ has finite size $k_m$. 

For the construction, we start with each $L_m = \emptyset$.  At any stage $s+1$,
when $(e,m,x)$ enters $R$, we add $x$ to the end of $L_m$. 

The argument for $\omega^*$ and $\zeta$ is similar. 

\end{proof}

\end{document}
